% Options for packages loaded elsewhere
% Options for packages loaded elsewhere
\PassOptionsToPackage{unicode}{hyperref}
\PassOptionsToPackage{hyphens}{url}
\PassOptionsToPackage{dvipsnames,svgnames,x11names}{xcolor}
%
\documentclass[
  letterpaper,
  DIV=11,
  numbers=noendperiod]{scrartcl}
\usepackage{xcolor}
\usepackage[top=30mm,left=30mm,right=30mm,bottom=30mm]{geometry}
\usepackage{amsmath,amssymb}
\setcounter{secnumdepth}{5}
\usepackage{iftex}
\ifPDFTeX
  \usepackage[T1]{fontenc}
  \usepackage[utf8]{inputenc}
  \usepackage{textcomp} % provide euro and other symbols
\else % if luatex or xetex
  \usepackage{unicode-math} % this also loads fontspec
  \defaultfontfeatures{Scale=MatchLowercase}
  \defaultfontfeatures[\rmfamily]{Ligatures=TeX,Scale=1}
\fi
\usepackage{lmodern}
\ifPDFTeX\else
  % xetex/luatex font selection
\fi
% Use upquote if available, for straight quotes in verbatim environments
\IfFileExists{upquote.sty}{\usepackage{upquote}}{}
\IfFileExists{microtype.sty}{% use microtype if available
  \usepackage[]{microtype}
  \UseMicrotypeSet[protrusion]{basicmath} % disable protrusion for tt fonts
}{}
\makeatletter
\@ifundefined{KOMAClassName}{% if non-KOMA class
  \IfFileExists{parskip.sty}{%
    \usepackage{parskip}
  }{% else
    \setlength{\parindent}{0pt}
    \setlength{\parskip}{6pt plus 2pt minus 1pt}}
}{% if KOMA class
  \KOMAoptions{parskip=half}}
\makeatother
% Make \paragraph and \subparagraph free-standing
\makeatletter
\ifx\paragraph\undefined\else
  \let\oldparagraph\paragraph
  \renewcommand{\paragraph}{
    \@ifstar
      \xxxParagraphStar
      \xxxParagraphNoStar
  }
  \newcommand{\xxxParagraphStar}[1]{\oldparagraph*{#1}\mbox{}}
  \newcommand{\xxxParagraphNoStar}[1]{\oldparagraph{#1}\mbox{}}
\fi
\ifx\subparagraph\undefined\else
  \let\oldsubparagraph\subparagraph
  \renewcommand{\subparagraph}{
    \@ifstar
      \xxxSubParagraphStar
      \xxxSubParagraphNoStar
  }
  \newcommand{\xxxSubParagraphStar}[1]{\oldsubparagraph*{#1}\mbox{}}
  \newcommand{\xxxSubParagraphNoStar}[1]{\oldsubparagraph{#1}\mbox{}}
\fi
\makeatother


\usepackage{longtable,booktabs,array}
\newcounter{none} % for unnumbered tables
\usepackage{calc} % for calculating minipage widths
% Correct order of tables after \paragraph or \subparagraph
\usepackage{etoolbox}
\makeatletter
\patchcmd\longtable{\par}{\if@noskipsec\mbox{}\fi\par}{}{}
\makeatother
% Allow footnotes in longtable head/foot
\IfFileExists{footnotehyper.sty}{\usepackage{footnotehyper}}{\usepackage{footnote}}
\makesavenoteenv{longtable}
\usepackage{graphicx}
\makeatletter
\newsavebox\pandoc@box
\newcommand*\pandocbounded[1]{% scales image to fit in text height/width
  \sbox\pandoc@box{#1}%
  \Gscale@div\@tempa{\textheight}{\dimexpr\ht\pandoc@box+\dp\pandoc@box\relax}%
  \Gscale@div\@tempb{\linewidth}{\wd\pandoc@box}%
  \ifdim\@tempb\p@<\@tempa\p@\let\@tempa\@tempb\fi% select the smaller of both
  \ifdim\@tempa\p@<\p@\scalebox{\@tempa}{\usebox\pandoc@box}%
  \else\usebox{\pandoc@box}%
  \fi%
}
% Set default figure placement to htbp
\def\fps@figure{htbp}
\makeatother





\setlength{\emergencystretch}{3em} % prevent overfull lines

\providecommand{\tightlist}{%
  \setlength{\itemsep}{0pt}\setlength{\parskip}{0pt}}



 


\KOMAoption{captions}{tableheading}
\makeatletter
\@ifpackageloaded{caption}{}{\usepackage{caption}}
\AtBeginDocument{%
\ifdefined\contentsname
  \renewcommand*\contentsname{Table of contents}
\else
  \newcommand\contentsname{Table of contents}
\fi
\ifdefined\listfigurename
  \renewcommand*\listfigurename{List of Figures}
\else
  \newcommand\listfigurename{List of Figures}
\fi
\ifdefined\listtablename
  \renewcommand*\listtablename{List of Tables}
\else
  \newcommand\listtablename{List of Tables}
\fi
\ifdefined\figurename
  \renewcommand*\figurename{Figure}
\else
  \newcommand\figurename{Figure}
\fi
\ifdefined\tablename
  \renewcommand*\tablename{Table}
\else
  \newcommand\tablename{Table}
\fi
}
\@ifpackageloaded{float}{}{\usepackage{float}}
\floatstyle{ruled}
\@ifundefined{c@chapter}{\newfloat{codelisting}{h}{lop}}{\newfloat{codelisting}{h}{lop}[chapter]}
\floatname{codelisting}{Listing}
\newcommand*\listoflistings{\listof{codelisting}{List of Listings}}
\makeatother
\makeatletter
\makeatother
\makeatletter
\@ifpackageloaded{caption}{}{\usepackage{caption}}
\@ifpackageloaded{subcaption}{}{\usepackage{subcaption}}
\makeatother
\usepackage{bookmark}
\IfFileExists{xurl.sty}{\usepackage{xurl}}{} % add URL line breaks if available
\urlstyle{same}
\hypersetup{
  pdftitle={AI-Driven Rehabilitation System for Arabic Pediatric Communication Disorders},
  pdfauthor={Computer Systems Engineering Student},
  colorlinks=true,
  linkcolor={blue},
  filecolor={Maroon},
  citecolor={Blue},
  urlcolor={Blue},
  pdfcreator={LaTeX via pandoc}}


\title{AI-Driven Rehabilitation System for Arabic Pediatric
Communication Disorders}
\usepackage{etoolbox}
\makeatletter
\providecommand{\subtitle}[1]{% add subtitle to \maketitle
  \apptocmd{\@title}{\par {\large #1 \par}}{}{}
}
\makeatother
\subtitle{Senior Graduation Project - Chapter 1}
\author{Computer Systems Engineering Student}
\date{2026-07-01}
\begin{document}
\maketitle

\renewcommand*\contentsname{Table of contents}
{
\setcounter{tocdepth}{3}
\tableofcontents
}

\section{Introduction}\label{introduction}

\subsection{Overview}\label{overview}

Communication is a fundamental human right and a critical vector for
cognitive, social, and emotional development. Pediatric communication
disorders present unique developmental hurdles that, if left
unaddressed, can lead to permanent academic and social marginalization.
This project explores the integration of computer vision, speech signal
processing, and generative artificial intelligence to deliver an
accessible, automated, and dialect-aware rehabilitation tool tailored
for Palestinian children.

\subsection{Problem Statement}\label{problem-statement}

In the Arab world, and specifically within Palestine, pediatric
communication disorders pose a severe and expanding public health and
educational crisis. Children suffering from disorders such as Childhood
Apraxia of Speech (CAS), verbal stuttering, phonological deletion, and
environmental language delays face profound barriers to social
integration. Left untreated, these early childhood speech impediments
transform into permanent cognitive and social liabilities.

The core of the problem lies in a compounding deficit of clinical,
geographical, and technological resources:

\begin{itemize}
\tightlist
\item
  \textbf{Severe Clinical Shortages:} There is an acute shortage of
  certified Speech-Language Pathologists (SLPs) relative to the
  vulnerable pediatric population.
\item
  \textbf{Geographical and Economic Barriers:} Existing speech therapy
  infrastructure in Palestine is heavily centralized within major urban
  centers. Families residing in rural or marginalized areas face
  prohibitive transportation costs and long clinic waiting lists,
  leaving thousands of rural children entirely unserved.
\item
  \textbf{Dialectal Technological Void:} Modern digital therapy
  applications and AI speech tools are built almost exclusively for
  Western languages or Modern Standard Arabic (MSA). They lack the
  phonetic modeling, acoustic tolerances, and vocabulary profiles
  required to process the distinct phonetics and localized variations of
  the Palestinian/Levantine dialect spoken by children.
\end{itemize}

Consequently, there is an urgent, unmet engineering need to design and
implement an automated, accessible, intelligent, and dialect-aware
system capable of screening, tracking, and assisting in the remote
rehabilitation of Arabic pediatric communication disorders.

\subsection{Local Statistics (Palestine
Context)}\label{local-statistics-palestine-context}

To establish the critical socio-economic necessity of this engineering
intervention, this project anchors its problem statement in empirical
data compiled by the Palestinian Central Bureau of Statistics (PCBS) and
recent regional epidemiological research:

\begin{itemize}
\tightlist
\item
  \textbf{High Prevalence Rates:} According to census data from the
  PCBS, communication disabilities represent the \textbf{second most
  prevalent childhood disability} in Palestine, accounting for
  \textbf{37.3\% of all pediatric disability cases}.
\item
  \textbf{Geographical Disparity:} The prevalence fluctuates based on
  geographical vulnerability, affecting \textbf{35.7\% of disabled
  children in the West Bank} and rising to \textbf{40.1\% within the
  Gaza Strip}.
\item
  \textbf{Regional Vulnerability (Jenin Governorate):} The Jenin
  Governorate---the home region of the Arab American University
  (AAUP)---historically records the highest overall disability rate in
  the West Bank at \textbf{4.1\% of the total population}, creating an
  immediate, hyper-local demand for targeted technological intervention.
\item
  \textbf{Socio-Economic Barriers:} Local case-control studies reveal
  that children residing in rural Palestinian communities face a
  drastically elevated risk of prolonged speech and language
  delays---with an \textbf{Odds Ratio (OR) of 4.4}---directly driven by
  low family income and restricted physical access to urban speech
  centers.
\item
  \textbf{Consanguinity and Environmental Trauma:} Congenital factors,
  heavily correlated with local consanguinity rates, account for
  \textbf{33.6\% of communication disabilities} in Palestine, while
  acute environmental stressors and systemic trauma heavily exacerbate
  fluency disorders (stuttering) and sudden speech regression in
  children.
\end{itemize}

\subsection{Project Objectives}\label{project-objectives}

The primary goal of this senior project is to develop a
multi-dimensional, AI-driven, and low-latency software platform to
screen, gamify, and rehabilitate Arabic communication disorders in
children. The specific engineering and clinical objectives are detailed
below:

\subsubsection{Objective 1: Implement Real-Time Computer Vision for
Articulatory Assessment (The ``Saif''
Module)}\label{objective-1-implement-real-time-computer-vision-for-articulatory-assessment-the-saif-module}

To integrate the \emph{MediaPipe Face Mesh} framework to extract 3D
facial landmarks on edge mobile devices. The system will mathematically
calculate the Lip Pucker Ratio (LPR) in real time to visually monitor
and correct jaw drops and lip-rounding mechanics for children with
Childhood Apraxia of Speech (CAS) under the \emph{Digital DTTC
Protocol}.

\subsubsection{Objective 2: Build an Acoustic Pipeline for Fluency and
Stuttering Analysis (The ``Layla''
Module)}\label{objective-2-build-an-acoustic-pipeline-for-fluency-and-stuttering-analysis-the-layla-module}

To develop an audio signal processing engine using the \emph{Librosa
Python Library} to analyze Root Mean Square (RMS) energy envelopes and
Zero Crossing Rates (ZCR). The application will automatically detect
speech disfluencies (repetitions and blocks) down to 200 milliseconds of
latency, while validating vocal chord relaxation via the \emph{Easy
Onsets} clinical technique.

\subsubsection{Objective 3: Deploy Language Models for Phonological
Deletion Mapping (The ``Khaled''
Module)}\label{objective-3-deploy-language-models-for-phonological-deletion-mapping-the-khaled-module}

To fine-tune a pre-trained Arabic \emph{Wav2Vec 2.0} or \emph{Hubert}
speech-to-text model to perform automated \emph{Forced Alignment} on
children's speech. The system will utilize \emph{Levenshtein Distance}
matrices at the phoneme level to detect syllable deletions and trigger a
cognitive \emph{Minimal Pairs} visual feedback mechanism.

\subsubsection{Objective 4: Engineer an Adaptive Generative AI
Storytelling Engine (The ``Sara''
Module)}\label{objective-4-engineer-an-adaptive-generative-ai-storytelling-engine-the-sara-module}

To architect a dynamic backend utilizing conversational LLM APIs
(\emph{GPT-4o}) paired with localized Text-to-Speech (TTS) systems. The
framework will generate personalized, interest-driven interactive
stories that adapt their prompt-difficulty level dynamically based on
the child's recorded word counts to combat environmental language
delays.

\subsubsection{Objective 5: Design a Distributed Low-Latency Remote
Tele-Therapy Architecture (The ``Adam''
Module)}\label{objective-5-design-a-distributed-low-latency-remote-tele-therapy-architecture-the-adam-module}

To construct a robust real-time communication platform utilizing the
\emph{WebRTC Protocol}. The system will deploy real-time \emph{Gaze
Tracking Algorithms} to evaluate joint attention metrics, transmitting
instant biometric alerts regarding child distraction to the clinician's
dashboard over an optimized \emph{WebRTC DataChannel} in under 100
milliseconds.




\end{document}
